\SourceLeaf{365}%
Fällen führt sie unmittelbarer und auf kürzerem Weg zum Ziel als die Analysis, besonders wenn diese nicht geschickt eingesetzt wird. Dann hat sie eine ihr eigene Eleganz. Darüber hinaus bleibt sie vor allem in formaler Hinsicht und beim frühen Studium der Jugend unentbehrlich. Sie beugt Einseitigkeit vor, schärft den Sinn für Strenge und Klarheit und verleiht den Einsichten eine Lebendigkeit und Unmittelbarkeit, die analytische Methoden weit weniger fördern und mitunter sogar gefährden. Aus diesen Gründen sieht man mit Freude, dass einige französische Geometer in den letzten Jahrzehnten begonnen haben, sich mit besonderer Sorgfalt dem Teil der Geometrie zu widmen, der die Beziehungen zwischen Punkten und Linien außerhalb einer gemeinsamen Ebene, die Beziehungen verschiedener Ebenen zueinander sowie Linien doppelter Krümmung und gekrümmte Flächen untersucht. Soweit sie dabei ausschließlich geometrische Methoden anwenden, pflegen sie dieses Gebiet als eigenes Fach unter dem Namen \textit{Géométrie descriptive}. An dem vorliegenden Werk über diese Wissenschaft müssen wir besonders die große Klarheit und Knappheit der Darstellung, den wohlgeordneten Übergang vom Leichteren zum Schwierigeren und den Reichtum an neuen Sichtweisen und gelungenen Ausführungen loben. Wir empfehlen deshalb sein Studium als kräftige Nahrung für den Geist, die zweifellos viel dazu beitragen kann, den echten geometrischen Geist zu beleben und zu erhalten, den man in der neueren Mathematik sonst manchmal vermisst. Neben der rein wissenschaftlichen Seite dieser Untersuchungen ist ihr vielfältiger Nutzen für die Künste zu berücksichtigen, die es mit räumlichen Verhältnissen zu tun haben, insbesondere für das Zeichnen, die Landvermessung, die Baukunst und den Festungsbau. Auch in dieser Hinsicht hat der Verfasser seine Schrift durch verschiedene Anwendungen interessanter zu machen gewusst, obwohl er sie meist eher angedeutet als tatsächlich ausgeführt hat.
